% Part: lambda-calculus
% Chapter: lambda-definability
% Section: introduction

\documentclass[../../../include/open-logic-section]{subfiles}

\begin{document}

\olfileid[hi]{lam}{ldf}{int}
\olsection{परिचय}

पहली दृष्टि में लैम्ब्डा कलन केवल ऐसे व्यंजकों का अत्यंत अमूर्त कलन है जो
फलनों और अन्य वस्तुओं पर उनके अनुप्रयोगों को निरूपित करते हैं। लैम्ब्डा
कलन के वाक्यविन्यास में ऐसा कुछ नहीं जो संकेत करे कि ये किसी विशेष प्रकार
की वस्तुओं के फलन हैं; विशेषतः वाक्यविन्यास में प्राकृतिक संख्याओं का कोई
उल्लेख नहीं है। उसकी मूल संक्रियाएँ---अनुप्रयोग और लैम्ब्डा अमूर्तन---किसी
भी फलन पर लागू होती हैं, केवल प्राकृतिक संख्याओं के फलनों पर नहीं।

फिर भी कुछ चतुराई से लैम्ब्डा कलन में अंकगणितीय फलन, अर्थात् प्राकृतिक
संख्याओं पर फलन, परिभाषित करना संभव है। इसके लिए प्रत्येक प्राकृतिक
संख्या~$n \in \Nat$ हेतु हम एक विशेष $\lambd$-पद~$\num n$ परिभाषित करते
हैं, जो~$n$ का \emph{चर्च अंक} है। (चर्च अंकों का नाम अलॉन्ज़ो चर्च के नाम
पर है।)

\begin{defn}
  यदि $n \in \Nat$, तो अनुरूप \emph{चर्च अंक} $\num{n}$, ~$n$ को निरूपित
  करता है:
  \[
    \num{n} \ident \lambd[fx][f^n(x)]
  \]
  यहाँ $f^n(x)$, ~$x$ पर $f$ को $n$ बार लागू करने के परिणाम को दर्शाता
  है। उदाहरणार्थ, $\num{0}$, ~$\lambd[fx][x]$ और $\num{3}$,
  ~$\lambd[fx][f(f(f\,x))]$ है।
\end{defn}
  
चर्च अंक $\num n$ को ऐसे लैम्ब्डा पद के रूप में कूटित किया जाता है जो दो
तर्क $f$ और $x$ स्वीकार करके $f^n(x)$ लौटाने वाले फलन को निरूपित करता है।
चर्च अंक स्पष्टतः प्रसामान्य रूप में हैं।

लैम्ब्डा कलन में प्राकृतिक संख्याओं का निरूपण निश्चय ही तभी उपयोगी है जब
हम उनसे संगणना कर सकें। लैम्ब्डा कलन में चर्च अंकों से संगणना करने का अर्थ
है किसी $\lambd$-पद~$F$ को ऐसे चर्च अंक पर लागू करना और संयुक्त पद~$F\,
\num n$ को प्रसामान्य रूप में अपचयित करना। यदि वह सदा प्रसामान्य रूप में
अपचयित होता है और प्रसामान्य रूप सदा कोई चर्च अंक~$\num m$ है, तो संगणना
का निर्गम संख्या~$m$ माना जा सकता है। तब~$F$ को किसी फलन $f\colon \Nat
\to \Nat$ को परिभाषित करने वाला मान सकते हैं, अर्थात् ऐसा फलन कि $f(n) =
m$ तभी और केवल तभी जब $F\, \num n \red \num m$। चर्च--रॉसर गुण के कारण
प्रसामान्य रूप, यदि अस्तित्व में हों, अद्वितीय होते हैं। अतः यदि $F\, \num
n \red \num m$, तो कोई दूसरा प्रसामान्य पद, विशेषतः कोई दूसरा चर्च अंक,
ऐसा नहीं हो सकता जिसमें $F \, \num n$ अपचयित होता हो।

उलटे, कोई फलन $f\colon \Nat \to \Nat$ दिए होने पर हम पूछ सकते हैं कि क्या
कोई ऐसा पद $F$ है जो इस प्रकार~$f$ को परिभाषित करता है। उस स्थिति में हम
कहते हैं कि $F$, ~$f$ को \emph{!!{lambda define}} करता है और $f$,
!!{lambda definable} है। इसे बहु-स्थानी और आंशिक फलनों के लिए व्यापक किया
जा सकता है।

\begin{defn}
मान लें $f\colon \Nat^k \to \Nat$। हम कहते हैं कि कोई लैम्ब्डा पद~$F$,
$f$ को \emph{!!{lambda define}} करता है यदि सभी $n_0$, \dots,~$n_{k-1}$
के लिए,
\[
F \, \num{n_0} \, \num{n_1} \dots \num{n_{k-1}} \red \num{f(n_0, n_1, \dots,
  n_{k-1})}
\]
यदि $f(n_0, \dots, n_{k-1})$ परिभाषित है; और अन्यथा $F \, \num{n_0}
\, \num{n_1} \dots \num{n_{k-1}}$ का कोई प्रसामान्य रूप न हो।
\end{defn}

अत्यंत सरल उदाहरण अचर फलनों का है। पद $C_k \ident \lambd[x][\num{k}]$ उस
फलन $c_k\colon \Nat \to \Nat$ को !!{lambda define} करता है जिसके लिए
$c(n) = k$; क्योंकि किसी भी~$n$ के लिए $C_k \, \num n \ident
(\lambd[x][\num{k}])\num n \redone \num{k}$। तादात्म्य फलन
$\lambd[x][x]$ द्वारा !!{lambda defined} है। अधिक जटिल फलनों को परिभाषित
करना निश्चय ही कठिन है और प्रायः बहुत चतुराई माँगता है। इसलिए यह कुछ
आश्चर्यजनक है कि प्रत्येक संगणनीय फलन !!{lambda definable} है। इसका विलोम
भी सत्य है: यदि कोई फलन !!{lambda definable} है, तो वह संगणनीय है।

\end{document}
